\chapter{Anxiety}

\section*{Definition}
Anxiety is worry, fear, or unease about a possible threat or uncertainty. It is a normal response when proportionate to a situation, but an anxiety disorder is persistent, difficult to control, disproportionate, or impairing. Generalized anxiety disorder, panic disorder, phobias, social anxiety, obsessive-compulsive disorder, and trauma-related disorders are distinct diagnoses.

\section*{When to See a Doctor}

\begin{itemize}
 \item Worry or fear that persists, is difficult to control, or interferes with daily life
 \item Sleep disturbance, panic attacks, or physical symptoms such as rapid breathing, palpitations, muscle tension, gastrointestinal upset, tremor, or tingling
 \item Impairment at work, school, in relationships, or in self-care
 \item Safety checking behaviors reassurance seeking avoidance patterns escalating
 \item Suicidal ideation hopelessness accompanying severe depression anxiety comorbidity
 \item Panic attack frequency increasing or severity intensifying requiring emergency evaluation
\end{itemize}

\section*{How It Develops}
Anxiety involves interacting brain systems for threat detection, attention, learning, arousal, and emotional regulation. Stress activates autonomic and hormonal responses that can cause a racing heart, sweating, tremor, rapid breathing, and gastrointestinal symptoms. Genetics, learning, trauma, sleep, medical illness, substances, and context all contribute; no single neurotransmitter or brain pathway explains every anxiety disorder.

\section*{Causes}
\begin{itemize}
 \item Anxiety disorders and related conditions: generalized anxiety disorder, panic disorder, agoraphobia, specific phobias, social anxiety disorder, separation anxiety disorder, obsessive-compulsive disorder, PTSD, acute stress disorder, and adjustment disorder
 \item Medical causes: thyroid disease, low blood sugar, anemia, arrhythmias, asthma or other breathing disorders, menopause-related symptoms, pain, and medication effects
 \item Substance-related causes: alcohol or sedative withdrawal, stimulants, high caffeine intake, cannabis or other drugs, and medication withdrawal
 \item Genetic predisposition family history anxiety disorders twin studies show moderate heritability
 \item Environmental factors childhood trauma chronic stress adverse life events
\end{itemize}

\section*{Related Symptoms}
\begin{itemize}
 \item Restlessness fatigue concentration difficulties irritability muscle tension sleep disturbance palpitations sweating tremor tingling or numbness abdominal discomfort urinary frequency fear of impending doom dread anticipation negative rumination catastrophizing hypervigilance startle response
\end{itemize}
\textbf{Important:} Assessment should distinguish an anxiety disorder from medical or substance-related causes. Clinicians may use validated scales such as GAD-7, review medications and substances, and order targeted tests when indicated. CBT and medication options such as SSRIs or SNRIs are commonly used; benzodiazepines require particular caution because of sedation, dependence, and withdrawal risks. Urgent psychiatric assessment is needed for suicidality, psychosis, or inability to function safely.

\section*{Medical Evaluation}
\begin{itemize}
 \item Comprehensive psychiatric assessment includes history of present illness, past psychiatric history, substance use, medical history, and collateral information from family.
 \item Standardized screening tools (GAD-7, PHQ-9, PCL-5) help quantify symptom severity and track treatment response.
 \item Targeted medical evaluation may include thyroid testing, glucose, blood count, or an ECG when symptoms suggest a medical cause; toxicology testing is used selectively and does not by itself establish a diagnosis.
 \item Risk assessment for self-harm, suicide, and danger to others is an essential component of every psychiatric evaluation.
\end{itemize}

\section*{Treatment Options}
\begin{itemize}
 \item Cognitive-behavioral therapy is a first-line evidence-based treatment option for many anxiety disorders; other psychotherapies may be selected according to the diagnosis and patient preference.
 \item SSRIs or SNRIs are commonly used when medication is indicated. Other medicines are diagnosis-specific; benzodiazepines are generally short-term or exceptional options because of sedation, dependence, withdrawal, falls, and overdose risk, especially with opioids or alcohol.
 \item Combined treatment may help some people, but the best plan depends on severity, diagnosis, preference, access, and response; it is not automatically superior for everyone.
 \item Referral to psychiatry is indicated for severe symptoms, treatment resistance, suicidal ideation, or diagnostic uncertainty.
\end{itemize}

\section*{Self-Care and Home Management}
\begin{itemize}
 \item Establish a consistent daily routine with regular sleep, meals, and physical activity to support mental health stability.
 \item Practice stress management techniques including mindfulness, deep breathing exercises, and progressive muscle relaxation.
 \item Maintain social connections and avoid isolation by reaching out to trusted friends, family, or support groups.
 \item Avoid or reduce alcohol and recreational drugs, and notice whether caffeine worsens symptoms; do not abruptly stop alcohol or dependence-forming medicines without medical advice.
 \item If experiencing suicidal thoughts or crisis, contact emergency services or a crisis helpline immediately.
\end{itemize}

\section*{References}
\begin{thebibliography}{99}
\bibitem{anxiety:ref1} Craske MG, Stein MB. ``Anxiety.'' \textit{Lancet}. 2016;388(10063):3048-3059.
\bibitem{anxiety:ref2} Bandelow B, Michaelis S, et al. ``Treatment of anxiety disorders.'' \textit{Dialogues Clin Neurosci}. 2017;19(2):93-107.
\bibitem{anxiety:ref3} Stein MB, Sareen J. ``Generalized anxiety disorder.'' \textit{N Engl J Med}. 2015;373(21):2059-2068.
\bibitem{anxiety:ref4} Sadock BJ, Sadock VA, et al. \textit{Kaplan \& Sadock\'s Comprehensive Textbook of Psychiatry}, 10th ed. Wolters Kluwer, 2017.
\bibitem{anxiety:ref5} American Psychiatric Association. ``Practice guideline for the treatment of patients with panic disorder.'' \textit{Am J Psychiatry}. 2010;167(Suppl):1-68.
\bibitem{anxiety:ref6} National Institute for Health and Care Excellence. Generalised anxiety disorder and panic disorder in adults: management (CG113). Updated 2020.
\bibitem{anxiety:ref7} National Institute of Mental Health. Generalized anxiety disorder. Revised 2025.
\end{thebibliography}
